برای شبیه‌سازی یک ماشین تورینگ $T$ با $k$ نوار توسط ماشین تورینگ $M$ با یک نوار دوطرفه، الفبای نوار ماشین $M$ را به گونه‌ای گسترش می‌دهیم که بتواند موقعیت هدها و محتویات تمام نوارها را روی یک نوار واحد نگهداری کند.

\begin{itemize}
    \item به ازای هر علامت $X$ در الفبای نوار $T$، یک علامت نشانه‌دار $\dot{X}$ در $M$ در نظر می‌گیریم که نشان می‌دهد هد نوار مربوطه در آن لحظه روی این خانه قرار دارد.
    \item محتویات $k$ نوار مختلف بر روی تنها نوار ماشین $M$ با استفاده از یک علامت جداکننده به شکل زیر کنار یکدیگر چیده می‌شوند:
    $$ \# \dot{a}_{11} a_{12} \dots \# a_{21} \dot{a}_{22} \dots \# \dots \# a_{k1} \dots \dot{a}_{kn} \# $$
\end{itemize}

برای شبیه‌سازی هر یک گام از محاسبات ماشین $T$ توسط ماشین $M$:
\begin{enumerate}
    \item ماشین $M$ از ابتدای نوار شروع به حرکت کرده و کل نوار را اسکن می‌کند تا موقعیت تمام $k$ هد را پیدا کند و نمادهای زیر آن‌ها را در حافظه حالت خود ذخیره نماید.
    \item پس از خواندن $k$ نماد، ماشین $M$ با استفاده از تابع انتقال ماشین $T$، وضعیت بعدی، نمادهای جدید برای جایگزینی، و جهت حرکت هرکدام از هدها را تعیین می‌کند.
    \item ماشین $M$ مجدداً روی نوار حرکت کرده و نمادهای جدید را می‌نویسد، سپس موقعیت هدها را یک خانه به چپ یا راست جابجا می‌کند. اگر هدی به فضای خالی جدیدی نیاز داشته باشد، $M$ محتویات نوار را شیفت داده و فضا را باز می‌کند.
\end{enumerate}

فرض کنید ماشین تورینگ $T$ محاسبات خود را در $n$ گام انجام دهد:
\begin{itemize}
    \item در هر گام از اجرای $T$، ماشین $M$ برای پیدا کردن و به‌روزرسانی هدها نیازمند اسکن کردن کل طول نوار فعال خود است.
    \item طول کل نوار مورد استفاده در $M$ در بدترین حالت متناسب با تعداد کل گام‌های اجراست، یعنی حداکثر در مرتبه $O(n)$.
    \item بنابراین، برای شبیه‌سازی هر گام از اجرای $T$، ماشین $M$ به زمانی معادل $O(n)$ نیاز دارد.
    \item در نتیجه، کل زمان لازم برای اجرای $n$ گام توسط ماشین تک‌نواره $M$ برابر خواهد بود با:
    $$ O(n \cdot n) = O(n^2) $$
\end{itemize}
بنابراین، سربار زمانی روش شبیه‌سازی از مرتبه درجه دوم یا $O(n^2)$ است.
